\documentclass[10pt,a4paper]{amsart}
\usepackage[margin=19mm]{geometry}
\usepackage{amsmath,amssymb,mathtools}
\usepackage[scr=rsfs,cal=euler]{mathalfa}
\usepackage{xfrac}
\usepackage[unicode,hidelinks,bookmarks=false]{hyperref}
\newcommand{\ie}{i.e.\ }
\newcommand{\cf}{cf.\ }
\newcommand{\resp}{respectively\ }
\newcommand{\Subsection}{\subsection}
\providecommand{\nobreakdash}{\nobreak\hbox{-}}
\setlength{\emergencystretch}{2em}
\multlinegap0pt
\usepackage{bm}
\usepackage{bbm}
\usepackage{dsfont}
\usepackage{xspace}
\usepackage{cancel}
\usepackage{mathtools}
\usepackage{amsthm}
\makeatletter
\@ifundefined{thm}{\newtheorem{thm}{Thm}}{}
\@ifundefined{lem}{\newtheorem{lem}[thm]{Lemma}}{}
\makeatother
\def\bv{{\boldsymbol{v}}}
\def\bx{{\boldsymbol{x}}}
\def\uu{\mathbb{U}}
\title[A robust obstruction to full strong pluripotency for wild bl]{A robust obstruction to full strong pluripotency for wild blender-horseshoes}
\author{Shin Kiriki, Xiaolong Li, Yushi Nakano, Teruhiko Soma}
\date{Source: arXiv:2604.24063v1 (CC BY 4.0)}
\begin{document}
\maketitle

\noindent\emph{Source and licence.} A robust obstruction to full strong pluripotency for wild blender-horseshoes. Version arXiv:2604.24063v1 (CC BY 4.0), distributed under the Creative Commons Attribution 4.0 International licence, as recorded in the arXiv licence metadata for this exact version and shown on the abstract page https://arxiv.org/abs/2604.24063v1. Full rights notice: licenses/arxiv-2604.24063-CC-BY-4.0.txt.\par\smallskip

\noindent\textit{Editorial scope.} Counted unit: the Lem whose source label is \texttt{l\_ea\_horizontal} in the article (source0.tex lines 758--763, source label \texttt{l\_ea\_horizontal}), delivered together with the article's own complete original proof of it (source0.tex lines 764--866, one \texttt{proof} environment of 103 lines). No mathematics is written, repaired or re-derived: statement and proof are the article's own text line for line. Required context retained uncounted: eqn\_eigen\_v (lines 462--465); eqn\_f0V (lines 488--494); eqn\_min (lines 518--525); eqn\_d\_eta (lines 694--696); eqn\_phi\_2 (lines 699--702). Every definition, display and label of those lines is retained unchanged. Omitted from this package and not redistributed: 1--461, 466--487, 495--517, 526--693, 697--698, 703--757, 867--1176. Nothing inside a retained excerpt is omitted or altered: every retained body is delivered exactly as the article's lines are written, so no omission receipt of any kind accompanies this package. Citations: the retained excerpts carry no citation command at all, so no bibliography entry of the article is transcribed and no reference is redistributed. Reference adaptation: none was needed and none was made; every reference inside the delivered unit resolves inside it, so no unresolved reference or citation is produced. Printed slips kept literal: no printed word, symbol, hypothesis or number of the counted statement or of its proof was altered. Layout and class: the article's own class is \texttt{amsart} with packages including amsmath, amsthm, amscd, amssymb, mathrsfs, graphicx, hyperref, xcolor, enumerate, yhmath, color, lineno, bm; the wrapper is the lane's pinned \texttt{amsart} editorial wrapper at 10pt on A4, which restates the theorem-like environments the retained text needs and adds the article's own macro definitions that the retained text needs (\texttt{\textbackslash bv}, \texttt{\textbackslash bx}, \texttt{\textbackslash uu}), with extra wrapper packages (bm, bbm, dsfont, xspace, cancel, mathtools). The delivered numbering is the unit's own. Assets: no retained span contains a figure environment, a graphics-inclusion command, a PDF-page inclusion, a table or a TikZ picture; no image file is copied into this package, the package declares no asset dependency and no image file is redistributed with it. Licence: version arXiv:2604.24063v1 of this article is distributed under the Creative Commons Attribution 4.0 International licence, as recorded in the arXiv licence metadata for this exact version and shown on the abstract page https://arxiv.org/abs/2604.24063v1. A full rights notice is delivered at licenses/arxiv-2604.24063-CC-BY-4.0.txt. Characters: every retained excerpt is pure ASCII, so the delivered engine, XeLaTeX with its default Latin Modern text font, reports no missing character.
\begin{equation}\label{eqn_eigen_v}
\lambda_{\rm ss}<\frac12,\quad
\lambda_{\rm ss}<\lambda_{\rm cs0}\leq \lambda_{\rm cs1}<1<\lambda_{\rm cs0}+\lambda_{\rm cs1},\quad 2<\lambda_{\rm u}.
\end{equation}

\begin{equation}\label{eqn_f0V}
f_{0}(x,y,z)=
\begin{cases}
(\lambda_{\mathrm{u}}x,\lambda_{\mathrm{ss}}y,\zeta_0(z))&\text{if}\ (x,y,z)\in \mathbb{V}_{0},\\
(\lambda_{\mathrm{u}}(1-x),1-\lambda_{\mathrm{ss}}y,\zeta_1(z)) &\text{if}\ (x,y,z)\in \mathbb{V}_{1},
\end{cases}
\end{equation} 

\begin{align}
&a_{1}>0,\quad a_2>0,\quad a_3a_4<0,\label{eqn_aaa}\\
&|a_{3}|\left(\frac12+\varepsilon_0\right)<\frac12-\lambda_{\rm ss}(1+\varepsilon_0),\label{eqn_|a3|}\\
&|a_3|<\min\left\{\dfrac{a_2|a_4|}{\sqrt{4a_1^2+a_2^2}},\dfrac{a_1a_2}{\sqrt{4a_1^2+a_2^2}}(1-2\varepsilon_0)\right\},\label{eqn_min}\\
&a_1^{-1}a_2(1+\varepsilon_0+\mu)^2<\left(\frac12-\lambda_{\rm u}^{-1}\right)^2,\  
a_1^{-1}a_2a_4^2(1+\varepsilon_0+\mu)<\left(\frac12+\varepsilon_0\right)^2,\label{eqn_a42}\\
&0<-\varepsilon_0+\mu,\quad a_2(1+\varepsilon_0+\mu)<\lambda_{\rm u}^{-1}.\label{eqn_e_0mu}
\end{align}

\begin{equation}\label{eqn_d_eta}
\left|\frac{d\eta}{dx}(x)\right|\leq 2a_1\left(\frac12+\varepsilon_0\right)
\end{equation}

\begin{equation}\label{eqn_phi_2}
\varphi_{2}(x,y,z)=\left(\eta(x)+a_{2}(z+\mu),\ a_{3}\Bigl(y-\frac{1}{2}\Bigr)+\frac{1}{2},\ 
a_{4}\Bigl(x-\frac{1}{2}\Bigr)+\frac{1}{2} \right).
\end{equation}

\begin{lem}\label{l_ea_horizontal}
There exists a neighborhood $\mathcal{U}_0$ of $f_0$ such that, 
for any $\varepsilon$-almost horizontal proper surface $S$ in $\mathbb{R}^3$, $g\in\mathcal{U}_0$ and $i=0,1,2$, 
the surface 
$\varphi_{g,i}(S)$ is also $\varepsilon$-almost horizontal and proper.
\end{lem}

\begin{proof}
First we consider the case of $i=0,1$.
By \eqref{eqn_f0V}, one can choose $\mathcal{U}_0$ so that 
\begin{equation}\label{eqn_Dphi_i}
D\varphi_{g,i}(\bx)=
\begin{pmatrix} (-1)^i\lambda_{\rm u}+\alpha_{11}(\bx)& \alpha_{12}(\bx)& \alpha_{13}(\bx)\\
\alpha_{21}(\bx)&(-1)^i \lambda_{\rm ss}+\alpha_{22}(\bx)& \alpha_{23}(\bx)\\
\alpha_{31}(\bx)& \alpha_{32}(\bx) &\lambda_{{\rm cs}i}+\alpha_{33}(\bx)
\end{pmatrix}
\end{equation}
holds for any $g\in\mathcal{U}_0$ and $\bx=(x,y,z)\in \mathbb{R}^3$,  
where each $\alpha_{kl}:\mathbb{R}^3\longrightarrow \mathbb{R}$ is a continuous function supported on $\mathbb{U}_i$ and satisfying  
$\sup\bigl\{|\alpha_{kl}(\bx)|\, ;\, \bx \in \mathbb{R}^3\bigr\}=O(\varepsilon^2)$.
Here we denote by $O(\varepsilon^2)$ a non-negative constant with $\limsup\limits_{\varepsilon\to +0}
O(\varepsilon^2)/\varepsilon^2<\infty$.
The definition of $\boldsymbol{C}_{\varepsilon}^{\mathrm{uc}}(\boldsymbol{x})$ 
implies $|v^{\mathrm{s}}|\leq \varepsilon$ for any $\bv=(v^{\mathrm{u}},v^{\mathrm{s}},v^{\mathrm{cs}})\in 
\boldsymbol{C}_{\varepsilon}^{\mathrm{uc}}(\boldsymbol{x})$ with 
$(v^{\mathrm{u}})^2+(v^{\mathrm{cs}})^2=1$.
We set $D\varphi_{g,i}(\bx)(v^{\rm u},v^{\rm s},v^{\rm cs})^T
=(\widehat{v}^{\rm u},\widehat{v}^{\rm s},\widehat{v}^{\rm cs})^T$, 
where $(v_1,v_2,v_3)^T$ denotes the transposed column vector of a row vector $(v_1,v_2,v_3)$.
By \eqref{eqn_Dphi_i}, 
\[
|\widehat{v}^{\rm s}|\leq \lambda_{\rm ss}|v^{\rm s}|+O(\varepsilon^2)
\leq \varepsilon\lambda_{\rm ss}+O(\varepsilon^2).
\]
Moreover, 
\begin{align*}
\varepsilon\sqrt{(\widehat{v}^{\rm u})^2+(\widehat{v}^{\rm cs})^2}
&\geq 
\varepsilon\sqrt{\lambda_{\rm u}^2(v^{\rm u})^2+\lambda_{{\rm cs}i}^2(v^{\rm cs})^2}-O(\varepsilon^2)\\
&>\varepsilon\lambda_{{\rm cs}i}\sqrt{(v^{\rm u})^2+(v^{\rm cs})^2}-O(\varepsilon^2)\\
&
=\varepsilon\lambda_{{\rm cs}i}-O(\varepsilon^2).
\end{align*}
By \eqref{eqn_eigen_v},   
$|\widehat{v}^{\rm s}|\leq \varepsilon\sqrt{(\widehat{v}^{\rm u})^2+(\widehat{v}^{\rm cs})^2}$ holds 
if we take $\varepsilon$ sufficiently small.
It follows that $D\varphi_{g,i}(\bx)\bv\in \boldsymbol{C}_{\varepsilon}^{\mathrm{uc}}(\boldsymbol{x})$.




In the case of $i=2$, by \eqref{eqn_phi_2}, we have 
\begin{equation}\label{eqn_Dphi2}
D\varphi_{g,2}(\bx)=\begin{pmatrix} \dfrac{d\eta}{dx}(x)+\beta_{11}(\bx)&\beta_{12}(\bx)&a_2+\beta_{13}(\bx)\\
\beta_{21}(\bx)&a_3+\beta_{22}(\bx)&\beta_{23}(\bx)\\
a_4+\beta_{31}(\bx)&\beta_{32}(\bx)&\beta_{33}(\bx)\end{pmatrix},
\end{equation}
where each $\beta_{kl}:\mathbb{R}^3\longrightarrow \mathbb{R}$ is a continuous function supported on $\mathbb{U}_2$ and satisfying  
$\sup\bigl\{|\beta_{kl}(\bx)|\, ;\, \bx \in \mathbb{R}^3\bigr\}=O(\varepsilon^2)$.
We set again $D\varphi_{g,2}(\bx)(v^{\rm u},v^{\rm s},v^{\rm cs})^T
=(\widehat{v}^{\rm u},\widehat{v}^{\rm s},\widehat{v}^{\rm cs})^T$.
As in the previous case, 
\begin{equation}\label{eqn_vs_a3}
|\widehat{v}^{\rm s}|\leq \varepsilon |a_3|+O(\varepsilon^2)
\end{equation}
for any $\bv=(v^{\mathrm{u}},v^{\mathrm{s}},v^{\mathrm{cs}})\in 
\boldsymbol{C}_{\varepsilon}^{\mathrm{uc}}(\boldsymbol{x})$ with 
$(v^{\mathrm{u}})^2+(v^{\mathrm{cs}})^2=1$.

First we consider the case when $|v^{\rm u}|\geq \dfrac{a_2}{2a_1}|v^{\rm cs}|$ or equivalently 
$|v^{\rm u}|\geq \dfrac{a_2}{\sqrt{4a_1^2+a_2^2}}$.
By \eqref{eqn_min} and \eqref{eqn_Dphi2}, it holds that 
\[
|\widehat{v}^{\rm cs}|\geq |a_4v^{\rm u}|-O(\varepsilon^2)\geq \dfrac{a_2|a_4|}{\sqrt{4a_1^2+a_2^2}}-O(\varepsilon^2)>|a_3|
\]
if $\varepsilon$ is taken sufficiently small.
By \eqref{eqn_vs_a3}, we may also assume that 
\begin{equation}\label{eqn_vcs>v3}
\varepsilon |\widehat{v}^{\rm cs}|>|\widehat{v}^{\rm s}|.
\end{equation}



Next we suppose that $|v^{\rm u}|<\dfrac{a_2}{2a_1}|v^{\rm cs}|$ or equivalently 
$|v^{\rm cs}|> \dfrac{2a_1}{\sqrt{4a_1^2+a_2^2}}$.
Again by \eqref{eqn_min} and \eqref{eqn_d_eta}, \eqref{eqn_Dphi2}, 
\begin{align*}
|\widehat{v}^{\rm u}|&\geq a_2|v^{\rm cs}|-\left|\frac{d\varphi}{dx}(x)v^{\rm u}\right|-O(\varepsilon^2)
\geq a_2|v^{\rm cs}|-a_2\left(\frac12+\varepsilon_0\right)|v^{\rm cs}|-O(\varepsilon^2)\\
&> a_2\left(\frac12-\varepsilon_0\right)\dfrac{2a_1}{\sqrt{4a_1^2+a_2^2}}-O(\varepsilon^2)
=\dfrac{a_1a_2}{\sqrt{4a_1^2+a_2^2}}(1-2\varepsilon_0)-O(\varepsilon^2)\\
&>|a_3|-O(\varepsilon^2).
\end{align*}
Thus we have 
\begin{equation}\label{eqn_vu>v3}
\varepsilon |\widehat{v}^{\rm u}|>|\widehat{v}^{\rm s}|
\end{equation}
for any sufficiently small $\varepsilon$. 
By \eqref{eqn_vcs>v3} and \eqref{eqn_vu>v3}, 
$|\widehat{v}^{\rm s}|\leq \varepsilon\sqrt{(\widehat{v}^{\rm u})^2+(\widehat{v}^{\rm cs})^2}$.
It follows that  
$D\varphi_{g,2}(\bx)\bv\in \boldsymbol{C}_{\varepsilon}^{\mathrm{uc}}(\boldsymbol{x})$.


Since $\varphi_{g,i}=\varphi_i$ on $\mathbb{R}^3\setminus \uu_i$ for $i=0,1,2$, 
the restriction 
$\mathrm{proj}|_{\varphi_{g,i}(S)}:\varphi_{g,i}(S)\longrightarrow \mathbb{R}^2$ is proper and surjective.
Hence $\varphi_{g,i}(S)$ is an $\varepsilon$-almost horizontal proper surface.
This completes the proof.
\end{proof}

\end{document}
